\documentclass[12pt]{article}

% House style preamble
\input{.house-style/preamble.tex}

% Update PDF metadata
\hypersetup{
  pdftitle={Not every stable cluster is homeostatic: Stability, network order, and control in projectible kinds}
}

\title{Not every stable cluster is homeostatic:\\ Stability, network order, and control in projectible kinds}
\author{Brett Reynolds \orcidlink{0000-0003-0073-7195}%
\thanks{Contact: \href{mailto:brett.reynolds@humber.ca}{brett.reynolds@humber.ca}}\\
Humber Polytechnic \& University of Toronto}
\date{\today}

\begin{document}
\maketitle

\begin{abstract}
Homeostatic property cluster theory gave natural-kinds realism a powerful
anti-essentialist form, on which a kind can be real because its properties
cluster through mechanisms rather than because it has an essence. The word
\term{homeostatic} has tended to carry more
weight than the theory has earned, because it is used for several achievements
that are not the same: reliable co-occurrence, network order, maintenance, and
feedback-sensitive control. This paper treats projectibility as the target of kind
claims and offers a diagnostic ladder: identify the projection target; ask whether
a stable profile supports it; ask whether network order makes the profile
non-accidental; ask what, if anything, maintains the profile; and ask whether that
maintenance is corrective control. Homeostasis belongs only at the last step,
where perturbation-sensitive organization preserves a higher-scale relation through
lower-scale variation. The Watt governor calibrates control, glycemic regulation
tests strict biological control, interjections supply a negative linguistic case,
and agreement monitoring supplies a candidate positive one. The result is a
disciplined way of using the framework's insight without turning \term{homeostatic}
into a loose synonym for stable or projectible.
\end{abstract}

\noindent Keywords: natural kinds; homeostatic property clusters;
projectibility; homeostasis; causal control; linguistic categories

\section{Introduction}

The target of this paper is a quiet ambiguity in recent anti-essentialist kind
theory. Boyd's homeostatic property cluster framework, Slater's stable property
cluster proposal, and Khalidi's causal-network account all reject essences, and
they all aim to explain why some categories support projection while others do
not. Each, though, locates the basis for projection somewhere slightly different.

Boyd ties it to homeostatic mechanisms that maintain the cluster
\citep{boyd1991}; Slater deliberately sets the stability of a cluster apart from
the means by which it is maintained \citep{slater2015}; and Khalidi locates it in
the causal-network order among a kind's properties
\citep{khalidi2013,khalidi2018nodes}. The ambiguity lies less in these accounts
than in the shared vocabulary of \enquote{homeostatic property clusters}, which
covers all of these options at once and so leaves underdetermined what secures
projection in any given case.

The looseness is clearest in the homeostatic-mechanism component of the framework.
Boyd gives that component considerable flexibility, treating the clustering as
\enquote{metaphorically (sometimes literally)} homeostatic and allowing it to be
sustained by mutual property support, underlying mechanisms, or both
\citep{boyd1999}. Onishi and Serpico are therefore right that the notion of a
homeostatic mechanism in HPC theory has never been fully explicated
\citep{onishi2022}.

The claim is that not every stable cluster is homeostatic. The literature
need not explicitly identify every stable cluster with a homeostatic one for this
to matter. Shared HPC vocabulary still makes it too easy to slide from stable
profile to stabilizer, from stabilizer to homeostasis, and from homeostasis to
projectible kind, as though these were one achievement.

The order should run the other way. Start from a projection target, ask what
empirical profile supports it, ask what network order makes the profile
non-accidental, ask what maintains it, and only then ask whether the maintenance is
corrective control. This order keeps Boyd's projectibility frame in place while
asking that \term{homeostasis} be used more literally, in a way that makes a genuine
demand. Weinberger's recent account of homeostasis as a form of causal control is
what supplies that demand \citep{weinberger2026HomeostasisCausalControl}.

The contribution is therefore diagnostic rather than revisionary. It leaves HPC
theory in place and separates five questions that the shared vocabulary tends to
compress. What is the projection target? What stable profile supports it? What
network order, if any, explains why the profile is not an accidental list? What
maintains that profile across the relevant range of cases? Does the maintaining
process correct perturbations and thereby preserve a higher-scale relation? Only
this last question gives the word \term{homeostatic} anything to mean beyond
stability.

The paper's claim is not that HPC theory should become control theory. It is that
HPC theorists need a diagnostic for when control is actually present.

The scope is Boyd-style projectible realist kind claims. Many classifications earn
their keep by mapping variation, stabilizing measurement, organizing historical
comparison, or coordinating inquiry before they support strong projection. Those
uses are legitimate, but the diagnostic here concerns categories offered as real
kinds because they support prediction, explanation, induction, or controlled
intervention.

The cases that follow run across a machine, an organism, and two linguistic
categories. They are not offered as a survey of those domains. They are chosen
because they stress different rungs of the ladder and test whether a single,
domain-neutral diagnostic can keep projection, stable profile, network order,
maintenance, and corrective control apart. A distinction in the theory of kinds
ought to travel in this way, and if it held only in biology that would be a
limitation rather than a sign of rigour. The cases earn their place by what they
reveal about the diagnostic, not by representing their fields.

The order of what follows is straightforward. The next five sections develop
projectibility and the successive evidential demands, ending with corrective
control. A short section then states the diagnostic ladder before the remaining
sections work through the cases and ask what the distinction changes for HPC theory
more broadly.

\section{Projectibility first}

Boyd's framework works because projectibility does not require necessary and
sufficient conditions. A category can support induction when its associated
properties are connected non-accidentally, even where no single property is shared
by every member. On Boyd's accommodation thesis, successful classification is a
matter of fitting our inferential practices to causal structure, so that induction
and explanation \enquote{always require that we accommodate our categories to the
causal structure of the world} \citep[224]{boyd1991}. We construct the categories,
but whether they succeed depends on a structure we do not control, and it is that
structure which does the explanatory work.

The problem of \term{projectibility} comes from Goodman's new riddle of
induction, where a predicate can fit all observed cases and still fail to license
projection to unobserved ones \citep{goodman1983}.
Boyd's contribution was to ground projectibility in causal structure without
grounding it in essences, which is also what gave HPC theory its purchase against
the worry, which Boyd associated with Hacking, that where a kind's boundaries are
fuzzy its membership is fixed by \enquote{people in concert} rather than by nature
\citep{boyd1991}.

This is one reason it helps to begin from the projection target rather than from
mechanism. Mechanisms matter to the theory because they explain why projection is
reliable, so a category that supports no projection is not redeemed by having a
tidy causal story behind it, and a category that does support projection is
already doing much of the work a realist about kinds cares about. The task of the
rest of the paper is to ask, for categories that are clearly projectible, which
further evidential demands they meet.

The risk is that this question can be made too easy. If the only thing asked is
whether a cluster recurs, a great many categories will qualify, because recurrence
on its own is cheap. Stable co-occurrence may license some inference, but an
analysis still has to specify what is being projected from what before the
category becomes theoretically interesting, and that happens only when its profile
supports prediction, explanation, or controlled intervention.

\section{Observed stability and invariance}

Slater's stable property cluster account presses the useful point that homeostatic
mechanisms are not always needed for natural kindness. What matters on this view is
the stability of a property cluster, considered apart from the various ways that
stability might be maintained, since a cluster can, in Boyd's phrase, hold itself
together \citep{slater2015}. This yields the first distinction the paper needs,
because a stable and projectible cluster may be perfectly real without being
homeostatic in any stronger sense.

Once the target is fixed, the next question is whether a stable profile supports
it. Stability that is left unexplained is nonetheless too thin to do all the work,
and Onishi and Serpico turn this into a dilemma.

If the explanatory role of underlying processes is made optional or postponed, as
stable-cluster views invite us to do, the account risks leaving underexplained why
some inductions are rational rather than lucky. If causal structure is retained
through causal-network order, as Khalidi proposes, the analysis still faces
interest-relativity, because which structure counts as the relevant one depends on
the explanatory purpose at hand \citep{onishi2022}. Their broader observation is
that the notion of a homeostatic mechanism has never been fully explicated, and
that interest-relativity is unlikely to be eliminable.

The diagnostic offered here does not try to remove that interest-relativity so
much as to make it explicit. Asking first what a category lets us project, and
against which declared expectation, fixes the explanatory target, and the question
of securing structure is then asked relative to that target rather than in the
abstract. Interest-relativity becomes a parameter the analysis declares rather than
a cost it pays without noticing.

Woodward's account of invariance helps locate what mere stability leaves out. The
kind of stability that matters for explanation is invariance under intervention,
not simply persistence across space and time, so that a generalization earns its
explanatory keep when it would continue to hold across a relevant range of changes,
including changes that could be imposed \citep{woodward2001}. The first diagnostic
step after the projection target is correspondingly modest: a stable property
cluster should support some projection across a specified range of contexts and
interventions, with evidence that goes beyond recurrence and shows that observing
part of the profile licenses a useful expectation about another part, or about how
the system behaves under a relevant change.

\section{Network order without control}

Khalidi's network account addresses a different shortcoming. A cluster is more than
a list of associated properties, because its properties can be ordered, with some
generating, sustaining, constraining, or explaining others. This is an advance over
a bare cluster picture, since it explains why some properties are diagnostic and
others merely derivative \citep{khalidi2013,khalidi2018nodes}. Khalidi himself
replaced the notion of a homeostatic mechanism with causal-network order in
pursuit of a more general account, so identifying that order as one securing
structure among several is continuous with his aim.

That order is the strength of the account. The present point is narrower: ordered
dependence, even when objective and explanatory, does not by itself amount to
feedback-sensitive correction.

Ordered dependence is still some distance from corrective control. A node in a
causal network can support projection because the properties around it stand in
systematic causal relations, and that is compatible with the system never tracking
a departure, never allowing lower-scale variables to adjust, and never preserving a
higher-scale relation through any such adjustment. Network order can therefore
explain projectibility without amounting to homeostasis.

The network-order step accordingly asks whether the profile has internal structure,
and the relevant evidence consists of asymmetries among the properties: which are
upstream, which derivative, which central, and which only peripheral cues. The
characteristic failure is the use of causal vocabulary for what is really an
unordered list, since redescribing co-occurrence in causal terms does not by itself
establish that any ordering obtains.

I use \term{network order} broadly across the diagnostic. In strictly biological or
mechanical cases the ordering will usually be causal, whereas in linguistic and
social cases the relevant dependencies may instead be functional, conventional,
sequential, or evidential. What matters in each case is that the profile has a
structure that explains projection, because a merely associated bundle of features
is not enough.

The broad use still needs a constraint. A non-causal ordering relation should be
asymmetric: one feature conditions, licenses, constrains, or structures another.
Evidential usefulness alone is not network order. A cue is evidentially privileged
only when it tracks an ordering relation in the profile rather than merely
co-occurring with the rest of the bundle.

\section{From stabilizers to controllers}

It helps to separate two questions that talk of maintenance tends to run together.
Whatever keeps a cluster clustered can be called its \term{stabilizer}, in the
sense that removing it would let the category fragment, drift, or dissolve. The
stabilizer is a functional role, and the causal structure that fills the role is
the mechanism. This is the functional-role description of what Boyd called a
homeostatic mechanism \citep{boyd1991}, and it has become a familiar way of talking
about how linguistic and social categories are held together. Maintenance, in this
sense, is the umbrella term.

The question the paper presses is which stabilizers are also controllers. A
\term{controller} is the narrower case: a stabilizer with feedback-sensitive
corrective organization that preserves a higher-scale relation through lower-scale
variation. Not every stabilizer meets this condition. Reliable co-occurrence can be
maintained by entrenchment, by transmission, or by a frozen accident that nothing
actively corrects, and ordered dependence can be maintained by causal wiring that
carries no error signal and returns the system to no target.

The two ideas answer different counterfactuals, which is why they come apart. The
stabilizer question is about the maintainer, and it asks what would happen to the
cluster if the maintaining process were removed. The controller question is about
the profile under threat, and it asks whether a perturbation to the profile is
sensed and answered in a way that brings the profile back toward a relation worth
projecting from.

A cluster can answer the first question and fail the second, since entrenchment and
transmission may hold a profile together without any of them detecting or correcting
a departure. The space between a profile that is merely held in place and one that
is actively corrected is where \term{homeostatic} tends to be overclaimed.

Read this way, maintenance is the fourth question in the ladder, not the first. A
stable profile can support projection before its maintainer is known, and a
network-ordered profile can explain why projection is reliable without yet showing
what keeps the network in place. Once a maintainer is identified, the final
question is whether that maintainer is a controller. Only the controller variety
earns \term{homeostatic} in anything like a literal sense, and the next section
says what control adds, using the account that states the standard most carefully.

\section{Corrective control}

Control-secured kinds are more demanding, and the first thing to get right is the
level at which the control is located. Kinds do not themselves sense, compare, or
correct anything. A kind earns the label only when its projectible profile is
maintained by a system or a practice that has feedback-sensitive corrective
organization, so that the bearer of the control is the system rather than the
category named.

Weinberger's central lesson is that homeostasis requires more than constancy,
because a controlled system can preserve a higher-scale relation precisely by
letting its lower-scale variables move. His dynamic causal models show that
feedback systems, and the Watt governor in particular, produce control relations
that static causal models and a simple vocabulary of intervention tend to obscure
\citep{weinberger2026HomeostasisCausalControl}.

Two consequences are worth drawing out. The first is that holding a variable
rigidly fixed can be a sign of less control rather than more, since appropriate
variation in the lower-scale variables is often the means by which the higher-scale
relation is kept in place. The second is that not every intervention is equally
probative. Perturbations to workload or demand can reveal the controller's response,
while clamping the regulated variable itself can sever the very feedback that
constitutes the control. Stillness under a clamp is therefore not evidence that
nothing was being controlled.

For the purposes of this paper, corrective control has four evidential marks. There
must be a perturbation that threatens the projectible profile; a variable or
relation that the system tracks and that matters to the profile; a response pathway
that changes lower-scale conditions; and a result in which that response preserves a
higher-scale relation supporting projection. Feedback on its own does not satisfy
these conditions, because what has to be shown is that the feedback explains why the
profile remains projectible when it is perturbed.

The biological control literature makes the same warning concrete. Control
mechanisms sense conditions, transmit signals, and alter constraints within the
mechanisms they control \citep{bichBechtel2022control}, and in organisms these
controls are integrated across the system rather than issued from a single command
point \citep{bichMossioEtAl2016,bichBechtel2022organization}. This is a good deal
more specific than the claim that a property cluster is stable, or that its
properties stand in some causal order, and it is the final and most demanding step
in the diagnostic ladder.

\section{The diagnostic ladder}

The proposal is best read as a diagnostic ladder rather than as an exhaustive
taxonomy. Control-secured cases will usually display stability, network order, and
maintenance as well, but the inference does not run in reverse. A stable cluster
need not have an internally ordered profile, a network-ordered kind need not have a
specified maintainer, and a maintained kind need not be maintained by corrective
control.

The ladder isolates five questions that are recurrently compressed in HPC
discussions: what is projected, what profile supports the projection, what order
makes the profile non-accidental, what maintains the profile, and whether that
maintenance is corrective control. Selection, drift, institutional scaffolding,
exemplar effects, and path dependence may each stabilize projectible profiles in
ways that deserve separate treatment. These cases matter because they may secure
projectibility without fitting neatly into the strict control role; the ladder is
meant to prevent conflation, not to classify every stabilizing process.

\begin{table}[t]
\centering
\small
\begin{tabular}{>{\raggedright\arraybackslash}p{0.21\linewidth}
                >{\raggedright\arraybackslash}p{0.35\linewidth}
                >{\raggedright\arraybackslash}p{0.30\linewidth}}
\toprule
Diagnostic step & Evidence needed & Failure mode \\
\midrule
Projection target & A specified inference, prediction, explanation, or
intervention the category is supposed to support & Kind talk without saying what
the category lets us project \\
Stable profile & Reliable co-occurrence across relevant contexts, with a
projection from observed to unobserved features & Mere recurrence without
projectible payoff \\
Network order & Asymmetries among properties, with some upstream, central,
derivative, or evidentially privileged & Causal vocabulary used for an unordered
list \\
Maintenance & A stabilizer whose removal would make the profile fragment, drift,
or dissolve & Treating persistence as if a maintainer had been specified \\
Corrective control & Perturbation-sensitive correction that preserves a
higher-scale relation through lower-scale variation & Treating stability,
maintenance, or feedback as control \\
\bottomrule
\end{tabular}
\caption{A diagnostic ladder for projectible profiles and homeostatic control.}
\label{tab:diagnostics}
\end{table}

The ladder is meant to turn the question around. For any proposed HPC kind, it asks
five things before the word \term{homeostatic} is used. The first is what the
category lets us project. The second is what stable profile supports that
projection. The third is whether the profile has ordering relations that explain
why the projection is not accidental. The fourth is what maintains the profile. The
fifth is whether the maintainer corrects perturbations. The aim is not to police a
word but to make the evidential demand behind it explicit.

Nothing in the argument requires every historical use of \term{homeostatic} in HPC
theory to have meant feedback control. The claim is normative and diagnostic:
corrective control is the structure that gives the term distinctive content once
stability, maintenance, and order have been separated.

\section{How to read the cases}

The cases that follow do not all have the same standing, and the differences are
part of the point. The Watt governor is not put forward as a kind at all; it serves
to calibrate the control structure that a homeostatic-control-secured kind would
have to involve. Glycemic regulation is the strict biological control case, though
one in which the familiar set-point picture is itself under pressure. Interjections
are a linguistic category whose stable projectibility and partial network order are
clear enough, and they test why those rungs do not yet earn the stronger label.
Agreement monitoring is a candidate positive linguistic case, where the stronger
label comes closest to being earned once the relevant repair loop is specified.

They are not all kinds in the same sense. They are arranged around the evidential
structure required for kind claims: whether a profile supports projection, what
order explains that projection, what maintains the profile, and whether the
maintainer is corrective control. The governor calibrates the final question; the
other cases test how far biological and linguistic profiles move up the ladder.

Running these cases across a machine, an organism, and two linguistic categories is
deliberate. The diagnostic concerns how projectibility is secured, and that is not a
question internal to any one science. If the same ladder discriminates among a
governor, a regulatory profile, and two kinds of linguistic patterning, then the
distinction is doing work in the theory of kinds rather than reporting the folklore
of a particular field.

\section{Calibration case: Watt governor}

The Watt governor is a useful calibration case because its control structure is
explicit. A change in engine speed moves the governor's flyballs; their position
changes the valve aperture; the aperture changes the flow of steam; and the steam
flow changes the speed of the shaft. The controlled target is shaft speed, mediated
by the admission of steam, and a clamp placed directly on shaft speed would sever
the incoming causal relations and destroy the control relation altogether
\citep{weinberger2026HomeostasisCausalControl}.

That point does not make the governor immune to intervention. A change in workload
or demand is the right kind of perturbation, because it lets the feedback relation
operate and reveals whether lower-scale adjustment preserves the higher-scale
coupling.

The governor shows clearly why \term{homeostatic} ought not to mean constant. Its
lower-scale variables are allowed to vary precisely so that workload and steam flow
remain coupled in a stable higher-scale relation. This is the pattern the diagnostic
uses to calibrate its final rung, in which a perturbation is met by responsive
adjustment in the lower-scale variables, and a relation that supports projection is
preserved through that adjustment.

Because the governor is an artifact with a target fixed by its designer, it shows
the control structure in an especially clean form, and that is what makes it a
reference point rather than a kind. A candidate kind can then be tested against it,
and the question for any proposed homeostatic kind becomes whether it exhibits this
structure, rather than whether its properties happen to stay put.

\section{Biological extension: glycemic regulation}

Glycemic regulation is included not to show that biological homeostatic cases are
easy, but to show why the control diagnostic should be relation-preserving rather
than set-point-preserving. The relevant candidate is not glucose as a substance,
nor a single blood value, but the normoglycemic regulatory profile of an organism.
The standard endocrine picture
already supplies a control architecture, with a device that sets a normal range, a
sensor that measures the regulated variable, a detector that compares the measured
value against the range and emits an error signal, a controller, and an effector
that moves the variable back toward the range \citep{bichMossioSoto2020}. Taken at
face value, this looks control-secured.

The projectible payoff is the organism's metabolic response across changing
conditions. Given the normoglycemic regulatory profile, one expects coordinated
changes in insulin, glucagon, hepatic output, and glucose uptake across feeding,
fasting, and activity, rather than an isolated reading that happens to fall within
a range.

The standard model already reserves the word \term{controller} for the component
that reads the error signal and drives the effector. The vocabulary of the
calibration case and of the biological case
coincides, and not by accident, since both are feedback architectures in which a
sensed departure drives a correcting response. The question is never whether the
word fits the diagram, but whether the organism realizes the architecture that the
diagram describes.

Two complications make that question sharper. The first concerns the set point,
since population correlations for glucose and calcium may fit an equilibrium model
better than they fit a defended set point \citep{fitzgeraldBean2018}. The second is
that \textcite{bichMossioSoto2020} argue the feedback-loop diagram flattens a
hierarchical organism into a chain of components, encourages a search for neatly
localized parts, and assumes a normal range without explaining what establishes it.
They move instead to organizational closure, on which glucose regulation is part of
the organism's self-maintenance rather than the tracking of an externally given
number. Read this way, the corrective structure is real, but the set point is the
element the cybernetic analogy gets wrong.

That move marks the real contrast with the governor. The governor's target is set
from outside by a designer, whereas the organism's range is constituted by what
keeps the organism alive. The higher-scale relation that glycemic control preserves
is therefore not a number but a relation between glucose availability and the
organism's continued functioning, held in place through wide variation in the
lower-scale variables across feeding, fasting, and activity. A biological example
counts as homeostatic-control-secured only when the evidence displays this
organized correction, and not when the variable merely stays within range.

Glycemic regulation is therefore the strict biological case in the paper: the stable
profile is maintained by organized responses to perturbation, and what the system
returns to is best understood as the relation rather than the reading.

\section{Contrast case: interjections}

English interjections are the negative linguistic case. The projection target is
pragmatic and interactional: identifying a form as an interjection licenses
expectations about how it will be packaged, placed, and taken up in use. The stable
profile is familiar, since prosodic isolation, low syntactic integration,
non-inflection, non-referentiality, and interactional force tend to travel together
\citep{ameka1992,norrick2009}. Observing some of these properties licenses
expectations about the others, which makes the profile projectible rather than a
mere list.

The payoff is pragmatic. Forms such as \mention{huh?}, \mention{mmhm}, and
\mention{oh} support expectations about sequential position and interactional
uptake, including repair initiation, continuation, and a change in epistemic state
\citep{dingemanse2024}, while forms such as \mention{wow}, \mention{oops},
\mention{hey}, \mention{well}, and \mention{uh} invite expectations about stance,
event management, attention, transition, or delay. \textcite{norrick2009} is helpful
here because he treats interjections as an open-ended class with recurrent
pragmatic-marker functions.

The contrast between \mention{oh} and \mention{uh} shows the kind of projection at
issue. In an appropriate sequential position, \mention{oh} leads us to expect a
change-of-state uptake: the speaker has just registered, noticed, or reclassified
something. \mention{uh}, by contrast, leads us to expect delay or floor-holding
rather than a completed uptake. These are genuine projections from form and
placement to interactional role, but they do not yet show a controller preserving
the category profile.

The list is deliberately heterogeneous, since several of these forms are analysed as
discourse markers, fillers, summonses, or response tokens rather than as central
interjections. This heterogeneity is part of the test, because the contrast case is
not a classical kind with settled boundaries but a projectible interactional profile
whose kindhood is itself partly at issue.

The likely stabilizers include conventionalization, turn-taking routines, prosodic
recognizability, acquisition, register, uptake, and metapragmatic learning, and
together these make the category stable and useful. They also give parts of the
profile some network order, since sequential position, prosodic packaging, and
syntactic non-integration are not wholly independent of one another. That gets
interjections several rungs up the ladder. It does not yet put them at the top.

The homeostasis question is more demanding. For interjections, the higher-scale
relation would be an interactional profile linking turn position, prosodic packaging,
low integration, and uptake. The lower-scale variables would include token form,
intonation, sequential placement, phonological reduction, and local conventional
meaning, and the relevant perturbations would include syntactic integration,
inflection, shifted prosody, misplaced sequential use, or failure of uptake. A
homeostatic-control claim would require evidence that perturbations of this sort
reliably trigger coordinated correction that preserves the higher-scale interactional
relation. Repair, failures of uptake, teaching, and metalinguistic policing may
sometimes furnish control-like evidence, but the evidence has to be produced, and
until it is, interjections belong under the stable-profile, network-order, or
maintenance rungs.

The case is useful because it shows how far stable pragmatic projectibility can go
before a control claim has been earned.

\section{Candidate positive case: agreement monitoring}

Agreement monitoring is the linguistic case where the stronger label comes closest
to being earned. In agreement systems, one expression fixes features that other
expressions must express. In Corbett's terms, the source is the controller, the
agreeing expressions are targets, and canonical agreement is the case in which the
features take matching rather than mismatching values \citep{corbett2006}. The
features may be number, person, gender, noun class, or case. Gender is therefore
one possible feature in the system, not the object of the case.

That terminology is useful but not decisive. Corbett's controller-target relation
identifies a dependency structure in agreement. By itself, it is not a
feedback-control architecture. The controller does not sense a mismatch, compare
target forms with a standard, or issue a corrective signal. Agreement morphology
establishes the relation whose preservation matters; production monitoring and
self-repair are where a control-theoretic claim has to look for feedback.

The projectible payoff is clear. Once the controller's features and the agreement
domain are known, target forms can be projected across the domain. Number agreement
lets us project verb form from subject number. Noun-class agreement lets us project
class-marking across verbs, adjectives, and demonstratives
\citep{ashton1947,corbett2006}. The dependency is also ordered rather than merely
stable, because the controller-target relation is asymmetric: the controller fixes
the relevant feature values, and targets express them.

The first three rungs are therefore straightforward: agreement has a projection
target, a stable profile, and an ordered dependency relation. The harder question is
whether any maintenance process for that relation is also corrective control.

The control question arises only when that dependency is embedded in a corrective
process. Agreement-attraction errors show the perturbation clearly: speakers can be
drawn toward a nearby mismatching noun and produce an agreement error
\citep{bockMiller1991}. In online production monitoring, a speaker produces target
forms, monitors the developing utterance, detects trouble, and repairs or recasts
when the form fails to preserve the intended agreement relation
\citep{levelt1983}. On that loop, the tracked relation is agreement within the
utterance, the lower-scale variables are particular target forms, and the response
pathway is repair or recasting.

The agreement-specific character of the repair comes from the projected dependency
itself. Once the controller's features and the domain are fixed, target forms are
expected to express those values, so repair restores the higher-scale relation that
made the profile projectible rather than only smoothing local fluency.

This keeps constraint and correction apart. A grammar may mark a mismatch as
ill-formed without thereby supplying a corrective pathway. Corrective control
requires the additional process by which a departure is detected and answered. The
four marks may be met in the online repair loop: the mismatch threatens the
projectible profile, agreement is the tracked relation, repair changes the
lower-scale forms, and the higher-scale relation is preserved through that change.

Acquisition and normative correction are different candidate loops, not parts of a
single system unless coordination is shown. In acquisition, learners converge on
agreement patterns over developmental time; in normative contexts, teachers,
editors, or interlocutors may reject or correct mismatches. These processes help
stabilize agreement and may provide further control-like evidence, but they operate
at different scales from online monitoring. The honest conclusion is conditional:
agreement monitoring is not trivially homeostatic, but it supplies a specified
linguistic loop rather than a loose analogy from grammatical dependency to feedback
control.

Grammaticality at large is a harder and more contested candidate. If correction,
repair, acquisition, transmission, and uptake together formed an architecture that
returned usage to a community-sensitive profile, grammaticality could turn out to be
control-secured, but whether grammaticality is a single kind, and which loops would
be doing the controlling, is just what is in dispute. The diagnostic leaves the
possibility open without granting it, and agreement monitoring shows how a positive
case would have to specify its loop.

\section{What the diagnostic changes for HPC theory}

The ladder matters for the wider HPC programme in linguistics and social ontology.
Linguistic categories can have projection targets, stable profiles, network order,
and maintainers without all of them being homeostatic in the strict
control-theoretic sense, and social practices can be stabilized by records, offices,
sanctions, training, and uptake without thereby constituting single homeostatic
control systems. The diagnostic asks each case to display how far up the ladder it
goes rather than inferring control from the mere fact of persistence.

A natural objection is that social correction and metalinguistic policing are
themselves forms of feedback, and sometimes they are. The concession strengthens the
diagnostic rather than weakening it, because it means that some linguistic and social
categories may indeed be control-secured, but only where correction, repair,
acquisition, transmission, and uptake combine into a coordinated architecture of the
kind that agreement monitoring may exhibit. The point is that the evidence has to
earn the upgrade case by case, rather than the label being granted on the strength
of stability alone.

The distinction also bears on the more formal side of HPC work. It is tempting to
operationalize homeostasis as a return-to-basin claim, asking whether a variable
comes back to a set point after it has been perturbed. The governor and glycemic
regulation both show why that test is too narrow, since what each system preserves is
a relation, between workload and steam flow in the one case and between glucose
availability and self-maintenance in the other, held in place through variation in
the lower-scale variables.

A test that measures only the return of a variable to a fixed value can miss
control that works precisely by allowing the variable to move. The basin is better
defined over the relation the system preserves than over any set of fixed values.

For HPC work in linguistics this implies a division of labour. Many of the
categories the framework would call homeostatic are, on inspection, stable
projectible profiles, maintained profiles, or network-ordered kinds, which are real
and worth a realist's defence without being control-secured. Reserving
\term{homeostatic} for the cases that display corrective control leaves the others
untouched, since they remain real and projectible, while showing the analyst what to
look for, and what would count as having found it, before the strongest description
is reached for. What philosophers of science gain is a domain-neutral test that
distinguishes projection target, stable profile, network order, maintenance, and
corrective control, and that says what evidence would move a case from one rung to
another.

\section{Conclusion}

Projectibility remains the point throughout. A stable cluster can underwrite
projection without being homeostatic: what makes it projectible is the reliability
of the profile for a specified inferential purpose, whereas what makes it
homeostatic is a further corrective organization that preserves the relevant
profile under perturbation. Drawn this way, the diagnostic protects Boyd's
anti-essentialist insight while keeping \term{homeostatic} from becoming a loose
synonym for stable, useful, maintained, or causally organized.

\section*{Acknowledgements}

Drafting and revision benefited from OpenAI Codex (GPT-5), GPT-5.5 Pro, and Claude
Opus 4.8. All content has been reviewed and revised by the author, who takes full
responsibility for the final text.

\printbibliography

\end{document}
